%%% SPMP - Organisation: Beginn
\begin{center}
  \makeatletter
  \textbf{
    \@title
  }
  \makeatother
  \\[1em]
  Informatik -- Eingebettete Systeme -- Hochschule Bremerhaven
\end{center}

\begin{table}[ht]
  \centering
  \begin{tabular}[t]{|p{15mm}|p{30mm}|p{30mm}|p{50mm}|r|}
    \hline
    \textbf{Version} & \textbf{Autor*inn*en} & \textbf{geändert} & \textbf{Beschreibung} & \textbf{fertig am}\\
    \hline
    0.1              & A. B.            &  Abschnitte 1 bis 4  & Erstellung des Dokumentes & 01.10.2019\\
    \hline
  \end{tabular}
  \caption{Versionshistorie}
\end{table}

\vspace*{4em}

Relevante Standards und Normen
\begin{itemize}
\item IEEE 16326-2009 -- ISO/IEC/IEEE International Standard --
  Systems and Software Engineering -- Life Cycle Processes -- Project
  Management

  \cite{ieee-16326-2009-project-management}
\item IEEE 16085-2006 -- ISO/IEC 16085:2006, Standard for Software
  Engineering -- Software Life Cycle Processes -- Risk Management

  \cite{ieee-16085-2006-risk-management}
\end{itemize}
\clearpage
%%% SPMP - Organisation: Ende

\chapter{Einleitung}
\label{spmp:einleitung}

\section{Motivation für das Dokument}
\hinweis{In diesem Dokument liegt das Konzept, wie Sie die
  Projektaufgabe angehen, Ihre Entscheidungen nachvollziehbar
  gestalten und Transparenz herstellen wollen. Dieser Hinweistext muss
  für die letzte Abgabe im Semester gestrichen sein.}

\textbf{
  \hinweis{
    Diese Hinweise sind grün und haben in der fertigen Spezifikation
    nichts zu suchen.}
  }

\section{Projektübersicht}
\label{spmp:projektuebersicht}

\hinweis{Diese Übersicht zeigt das Projekt ganz allgemein.}

\section{Projekterzeugnisse}
\label{spmp:projekterzeugnisse}

\hinweis{In diesem Abschnitt geben Sie die abzuliefernden Dokumente,
  Programme und Hardware an.}

\chapter{Projektorganisation}
\label{spmp:projektorganisation}

\section{System-Prozessmodell}
\label{spmp:system-prozessmodell}

\hinweis{Hiermit sind der Software- und der
  Hardware-Entwicklungsprozess gemeint. Sie schreiben hier, nach
  welchen Standards und Vorgehensmodellen Sie arbeiten. }

\section{Rollen und Verantwortlichkeiten}
\label{spmp:rollen-verantwortlichkeiten}

\section{Werkzeuge und Techniken}
\label{spmp:werkzeuge-techniken}

\section{Vorkehrungen zur Arbeitssicherheit}
%\label{spmp:vorkehrungen-zur-arbeitssicherheit}
 
\begin{itemize}
  \item Im Labor dürfen wir essen und trinken, essen aber lieber vor der Tür, 
    sonst weg von den Computern
  \item Bei Verletzung bescheid sagen, es wird im Verbandbuch eintragen und wir 
    bekommen alles, was wir brauchen, aus dem Verbandkasten
  \item Beim Feuer stehen Feuerhandschuhe und Feuerlöscher zur Verfügung.
  \item Es gibt 2 Notausgänge und im Falle von Notfall versammeln wir uns auf der Wiese.
  \item Beim Notfall, anderen helfen
    \end{itemize}


\hinweis{Hier stehen als Liste mit Erläuterungen Werkzeuge und wofür
  diese benutzt werden und insbesondere auch, wenn Werkzeugketten
  eingesetzt werden, werden diese hier beschrieben. }

\hinweis{Werkzeuge und Techniken schließt auch alles ein, was Sie
  unternehmen, um die Qualität Ihres Projektes und Produktes
  sicherzustellen.}

\chapter{Projektmanagementplan}
\label{spmp:projektmanagementplan}

\section{Aufgaben}
\label{spmp:aufgaben}

% Kopiere Aufgabenbeschreibung von Beginn bis Ende für jede Aufgabe
% Aufgabenbeschreibung: Beginn
\subsection{Aufgabe-N}
\label{spmp:aufgabe-n}

\subsubsection{Beschreibung}
\label{spmp:aufgabe-n-beschreibung}
\hinweis{Sie beschreiben hier die Aufgabenstellung. Die Ergebnisse
  fließen in den Rest des Dokuments.}

\subsubsection{Erzeugnisse und Meilensteine}
\label{spmp:aufgabe-n-erzeugnisse}
\hinweis{Erzeugnisse sind z.B. Dateien, zu erstellende
  Hardware-Bausteine usw. Meilensteine sind bestimmte Zeiten, zu denen
  Erzeugnisse und Zwischenergebnisse fertig sein müssen. Meistens
  werden Meilensteine nicht von jemand anderem festgelegt -- das sind
  die wenigsten -- sondern von Ihnen als Projektmitarbeiter*innen
  selbst, um Ihre Arbeit zu strukturieren. }

\subsubsection{Benötigte Ressourcen}
\label{spmp:aufgabe-n-ressourcen}

\subsubsection{Abhängigkeiten und Bedingungen}
\label{spmp:aufgabe-n-abhaengigkeiten}

\subsubsection{Risiken}
\label{spmp:aufgabe-n-risiken}
\hinweis{Risiken können vielfältig sein. Hat z.B. das Labor, in dem
  Sie arbeiten nur ein einziges bestimmtes Gerät, dann ist der Ausfall
  dieses Geräts ein Risiko. Wenn Sie Risiken nennen, dann benötigen
  Sie auch Gegenmaßnahmen, damit diese Risiken nicht eintreten. Ein
  Risiko könnte z.B. sein, dass Sie krank werden. Und eine
  Gegenmaßnahme wäre dann möglicherweise, dass Sie auf sich acht
  geben, früh schlafen gehen, Obst essen etc. }

% Aufgabenbeschreibung: Ende

\section{Zeitplan}
\label{spmp:zeitplan}
\hinweis{Der Zeitplan kann anfangs nur aus einem Anfangsdatum und
  einem Enddatum bestehen. Gehen Sie nach dem Wasserfall vor, dann
  lassen sich oben genannte Meilensteine im Zeitplan
  wiederfinden. Gehen Sie anderweitig vor, dann können das Termine zu
  Zwischenständen oder Zwischenbesprechungen sein.
}

\chapter{Zusatzmaterial}
\label{spmp:zusatzmaterial}

